\section{Exact two-class ambiguity and optimal recovery error}\label{sec:sharp-ambiguity}
Here the cross factor is independent and uniform simple. Fix $a>b>0$ and put
\begin{equation}\label{eq:ambiguity-exponent}
 D=a+b,\quad\delta=a-b,\quad t=\delta+1,\quad\nu=(D-2)t,\qquad
 C_{a,b}=\frac{((bt)!)^2(a!)^{2t}}{(t!)^2(b!)^{2t}a^{t\delta}}.
\end{equation}
Let $K(G)$ count valid unordered $(a,b)$ bisections, including the plant.

\begin{theorem}[Exact ambiguity asymptotic]\label{thm:sharp-ambiguity}
Along every admissible block-size sequence,
\begin{equation}\label{eq:sharp-ambiguity}
 \Prob\{K(G)\ge2\}\sim C_{a,b}N^{-\nu},\qquad
 \Prob\{K(G)\ge3\}=O_{a,b}(N^{-\nu-1}).
\end{equation}
Under a uniformly hidden balanced plant, the optimal exact-recovery
Bayes error satisfies
\[
 R_N\sim(C_{a,b}/2)N^{-\nu}.
\]
The full-list algorithm
of Theorem~\ref{thm:fast-full-list}, returning any listed partition,
attains $R_N$ exactly in expected $\widetilde O_{a,b}(n^2)$ bit work.
A deterministic worst-case polynomial estimator attains the same leading
Bayes constant. For $(a,b)=(2,1)$, ambiguity is asymptotic to $4/N^2$
and Bayes error to $2/N^2$. These constants use block size $N=n/2$.
\end{theorem}

\subsection{Minimum flips and the negligible remainder}
For a second sign vector $\tau$, put
$w=(\sigma-\tau)/2=\one_{S_+}-\one_{S_-}$. Both signs are balanced,
so $|S_+|=|S_-|=j$; replace $\tau$ by $-\tau$ to take $j\le N/2$.
At a supported plus vertex,
\[
 \delta=|N(v)\cap S_+|-|N(v)\cap S_-|\le j-1.
\]
Thus $j\ge t$. The support estimate in the proof of
Theorem~\ref{thm:boolean}, starting at $2t+2$ instead of $2t$, gives
$O(N^{-(D-2)(t+1)})$ for nonminimum supports up to a small fixed
linear size. Its size-$h$ summand is
$[A(h/N)^{(D-2)/2}]^h$, with a geometric consecutive-ratio bound
in that range. The Boolean entropy gap gives an exponentially small
macroscopic remainder. Since $D\ge3$, the probability of any
alternative with $j\ge t+1$ is $O(N^{-\nu-1})$.

At $j=t$, equality forces each $S_\pm$ to induce $K_t$ and forbids
edges between them. Every support vertex has exactly $b$ internal and
$b$ cross boundary edges. At an outside vertex $v$, the equation
$Aw=\delta w$ requires equal numbers $m_v$ of neighbors in its own-side
and opposite-side support sets. On each old side,
$\sum_{m_v>0}m_v=bt$. Call the flip leading if every positive $m_v=1$.
Its exterior sets $W_+,W_-$ each have size $bt$. Its complete incident
graph $H$ has
\begin{equation}\label{eq:motif-excess}
 v_0=2t+2bt,\qquad e_0=t\delta+4bt,\qquad e_0-v_0=\nu.
\end{equation}
It is connected and has minimum degree two: each exterior vertex joins
the two connected support cliques. Support vertices have motif degree
$D>2$ and exterior vertices degree two.

A nonleading minimum flip has at most $2bt-1$ exterior vertices but still
exactly $e_0$ prescribed edges. There are finitely many such simple typed
incidence patterns. Each has at most $v_0-1$ vertices; its vertex count
and finite prescribed-stub probability give expected count
$O(N^{-\nu-1})$ in the configuration law. Conditioning all three factors
simple costs a bounded factor. Impossible patterns contribute zero.
Thus all nonleading alternatives together have probability
$O(N^{-\nu-1})$.

\subsection{Count every leading flip set once}
Choose support sets $S_\pm$ of size $t$ and exterior sets $W_\pm$ of
size $bt$. Four functions specify the boundary endpoints:
\[
 W_+\to S_+\text{ internal},\quad W_+\to S_-\text{ cross},\quad
 W_-\to S_-\text{ internal},\quad W_-\to S_+\text{ cross}.
\]
Each has exactly $b$ preimages at every support vertex, giving
$(bt)!/(b!)^t$ choices. The exact number of data choices is
\begin{equation}\label{eq:leading-data}
 A_N=\binom Nt^2\binom{N-t}{bt}^2
       \left[\frac{(bt)!}{(b!)^t}\right]^4.
\end{equation}
A realized leading flip determines all this data uniquely. Conversely
the prescribed edges make the flip valid: all support stubs are saturated
and each exterior vertex exchanges one internal and one cross edge.
Extra edges among exterior vertices are harmless. For $N>2t$ the
opposite orientation has $N-t>t$ flips per side, so no orientation or
drawing-automorphism factor belongs in \eqref{eq:leading-data}.

For one data choice the number of distinguished-stub assignments is
$c_H=(a!b!)^{2t}(ab)^{2bt}$. Each internal factor has
$e_{\rm int}=t\delta/2+bt$ prescribed edges; the cross factor has $2bt$.
One physical assignment has exact configuration probability
\begin{equation}\label{eq:motif-probability}
 P_N^*=\prod_{j=0}^{e_{\rm int}-1}(aN-2j-1)^{-2}(bN)_{2bt}^{-1}.
\end{equation}
The assignments are disjoint, because the required distinct neighbors
exhaust every support stub. After prescribing one, residual factors are
uniform and independent. Support residual degrees are zero; exterior
vertices have degrees $a-1,b-1$; all others keep $a,b$.
Every prescribed edge touches a saturated support vertex, so residual
edges cannot duplicate it. Full simplicity is exactly residual simplicity.

The residual lists differ from the original lists at only $O_{a,b}(1)$
vertices. Their loop and double-edge factorial moments have the same
limits: finite modifications change means by $o(1)$, and overlapping
fixed defect patterns lose a power of $N$. The bounded-degree ordinary
and bipartite simplicity estimates in~\cite{JansonSimple} give
\[
 s_N,s_N^*\longrightarrow
 \exp[-(a^2-1)/2-(b-1)^2/2]>0,\qquad s_N^*/s_N=1+o(1).
\]
The residual law is identical up to labels for all leading data and
prescriptions, so this ratio is uniform. Internal residual totals are
$aN-tD$, even because $tD=\delta(\delta+1)+2b(\delta+1)$ is even;
cross totals agree. No parity restriction has been added.

Consequently one data pattern has simple-model probability
\begin{equation}\label{eq:motif-mean}
 c_HP_N^*s_N^*/s_N
 =(1+o(1))(a!b!)^{2t}a^{-t\delta}N^{-e_0}.
\end{equation}
The exterior factor $(ab)^{2bt}$ cancels the boundary degree powers
in the denominators. For $X$ counting leading flip \emph{sets},
multiplication by \eqref{eq:leading-data} gives
\begin{equation}\label{eq:leading-mean}
 \E X\sim C_{a,b}N^{-\nu}.
\end{equation}
The simplicity ratio tending to one is essential for this coefficient.

\subsection{Factorial second moment and the two-state event}
If a finite connected graph $H$ has minimum degree at least two,
every proper subgraph $F$ with nonempty vertex set satisfies
\begin{equation}\label{eq:proper-excess}
 e(F)-v(F)\le e(H)-v(H)-1.
\end{equation}
A spanning proper subgraph omits an edge. Otherwise remove $r>0$
vertices. If $E_R$ edges lie inside them and $B$ leave them, then
$B>0$ and $2E_R+B\ge2r$, so $E_R+B>r$; at least $r+1$ edges are
removed. Further deletion only lowers excess. Retained isolated vertices
are included in $v(F)$.

For two distinct leading data patterns, let $F$ contain every common
vertex and edge. Their union has excess $2\nu-[e(F)-v(F)]$.
Each fixed overlap type contributes at most
$O(N^{-2\nu+e(F)-v(F)})$, by its vertex count and finite prescribed-stub
count; incompatible unions contribute zero. There are finitely many
types. Disjoint copies contribute $O(N^{-2\nu})$; proper nonempty
overlaps give $O(N^{-\nu-1})$ by \eqref{eq:proper-excess}.
There is no identical-image term: motif-degree-$D$ vertices identify
the support, which uniquely determines the data of a realized flip.
Distinct data with the same support cannot co-occur. Thus
\[
 \E[X(X-1)]=O(N^{-\nu-1}),\quad
 \Prob(X\ge1)=\E X+O(N^{-\nu-1}),\quad
 \Prob(X\ge2)=O(N^{-\nu-1}).
\]
Outside the union of the nonminimum and nonleading-minimum exceptional
events, $K(G)=1+X$.
These bounds and \eqref{eq:leading-mean} prove \eqref{eq:sharp-ambiguity}.

\subsection{Uniform posterior and the exact Bayes constant}
Hide the balanced plant uniformly. For an observed compatible labeled
graph, every candidate has the same likelihood,
\[
 [G_a(N)^2B_b(N)]^{-1}.
\]
Given that partition the three-factor decomposition is unique; the counts
do not depend on its vertex labels. The posterior on valid unordered
partitions is therefore exactly uniform. Its optimal conditional error
is $1-1/K(G)$, attained by choosing any valid partition. Hence
\begin{equation}\label{eq:Bayes-sandwich}
 0\le R_N-\tfrac12\Prob\{K(G)\ge2\}
 \le\tfrac12\Prob\{K(G)\ge3\}=O(N^{-\nu-1}).
\end{equation}
The factor one-half comes from two equal posterior possibilities on the
leading ambiguous event. Relabeling invariance equates ambiguity under
the hidden prior and a fixed plant; no extra factor-two correction occurs.
A publicly supplied fixed plant would not support this statistical claim.

The always-correct full-list algorithm returns any valid partition in
expected nearly quadratic work and attains $R_N$ exactly. For a
deterministic polynomial truncation, use the binary algorithm of
Section~\ref{sec:enumeration} with threshold $n^{\nu+1}$.
Its bounded binary moment gives interruption $O(N^{-\nu-1})$ and
elementary work $O_{a,b}(n^{\nu+2}+n^3)$. On each completed run
choose any valid partition, with optimal conditional Bayes error.
Interruption adds at most $o(N^{-\nu})$ to risk, giving the same leading
constant for this deterministic estimator.

The motif supplies an extra real eigenvector, hence a real-collision
lower bound of this order. The matching asymptotic here concerns Boolean
partition ambiguity; it does not give a matching real-corank failure
asymptotic or a prescribed-cross statistical lower bound.

